\documentclass[a4paper]{article}
% Español (México) con XeLaTeX: fontspec para fuentes Unicode, polyglossia
% para la división silábica y los nombres del documento en la variante mexicana.
\usepackage{fontspec}
\usepackage{polyglossia}
\setdefaultlanguage[variant=mexican]{spanish}
% Los nombres chinos van como «pinyin (original)»: xeCJK da a los caracteres
% Han su propia fuente; sin ella XeLaTeX los omite con un aviso.
% FandolSong no cubre algunos caracteres de nombres (U+73FA); WenQuanYi Zen Hei sí.
\usepackage[AutoFallBack=true]{xeCJK}
\setCJKmainfont{FandolSong-Regular.otf}
\setCJKfallbackfamilyfont{\CJKrmdefault}{WenQuanYi Zen Hei}
% polyglossia no traduce los nombres de listings.
\AtBeginDocument{\providecommand{\lstlistingname}{}\renewcommand{\lstlistingname}{Listado}%
  \providecommand{\lstlistlistingname}{}\renewcommand{\lstlistlistingname}{Índice de listados}}
\usepackage{amsmath, amssymb}
\usepackage{graphicx}
\usepackage[margin=2.35cm]{geometry}
\usepackage[most]{tcolorbox}
\usepackage{booktabs}
\usepackage{tabularx}
\usepackage{longtable}
\usepackage{array}
\usepackage{float}
\usepackage{xcolor}
\usepackage{hyperref}
\usepackage{fancyhdr}
\setCJKmainfont[AutoFakeBold=2.4]{AR PL UMing CN}
\setCJKsansfont[AutoFakeBold=2.4]{Droid Sans Fallback}
\setCJKmonofont[AutoFakeBold=2.4]{Droid Sans Fallback}
\hypersetup{colorlinks=true, linkcolor=blue!55!black, urlcolor=blue!60!black}
\graphicspath{{./}{figures/}}
\setlength{\parskip}{0.55em}
\setlength{\parindent}{2em}
\setlength{\emergencystretch}{3em}
\renewcommand{\arraystretch}{1.18}
\newtcolorbox{knowledgebox}[1]{enhanced, colback=blue!5!white, colframe=blue!70!black, colbacktitle=blue!70!black, coltitle=white, fonttitle=\bfseries, title={#1}, sharp corners, breakable}
\newtcolorbox{importantbox}[1]{enhanced, colback=yellow!10!white, colframe=yellow!75!black, colbacktitle=yellow!75!black, coltitle=black, fonttitle=\bfseries, title={#1}, sharp corners, breakable}
\newtcolorbox{warningbox}[1]{enhanced, colback=red!5!white, colframe=red!70!black, colbacktitle=red!70!black, coltitle=white, fonttitle=\bfseries, title={#1}, sharp corners, breakable}
\pagestyle{fancy}
\fancyhf{}
\fancyfoot[C]{\thepage}
\fancyfoot[R]{\footnotesize\color{gray}\href{https://github.com/hqhq1025/ai-course-notes}{AI Course Notes}}
\renewcommand{\headrulewidth}{0pt}
\renewcommand{\footrulewidth}{0pt}
\newcommand{\videourl}{https://www.youtube.com/watch?v=zrvnoYYPaWQ}
\newcommand{\repourl}{https://github.com/hqhq1025/ai-course-notes}

\begin{document}
\begin{titlepage}
\centering
{\Large Notas didácticas de revisión técnica\par}
\vspace{1.0cm}
{\huge\bfseries Un recorrido abierto por los artículos de investigación: cuatro ejes de la historia de la IA}\par
\vspace{0.5cm}
{\Large Zhang Xiaojun conversa con Xie Qingchi (谢青池)\par}
\vspace{0.35cm}
{\large Elaborado a partir del video público de Zhang Xiaojun Podcast, la transcripción generada en GPU, la hoja de ruta de artículos que figura en la descripción del video y diagramas conceptuales propios\par}
\vspace{0.3cm}
{\large 2025-10-28\par}
\vspace{0.85cm}
\includegraphics[width=0.88\textwidth,height=0.42\textheight,keepaspectratio]{cover.jpg}\par
\vfill
\begin{tcolorbox}[width=0.92\textwidth, colback=black!2!white, colframe=black!55, sharp corners]
\textbf{Título del video}: 117. Un recorrido abierto por los artículos de investigación: la historia completa de los cambios en paradigmas de modelos, Infra y datos, lenguaje y multimodalidad\par
\textbf{Canal del video}: Zhang Xiaojun Podcast\par
\textbf{Duración del video}: 04:22:37\par
\textbf{Enlace del video}: \href{\videourl}{\nolinkurl{\videourl}}\par
\textbf{Nota sobre los materiales}: YouTube no tiene subtítulos disponibles; el texto se elaboró a partir de la transcripción con faster-whisper large-v3 en CUDA, la hoja de ruta de 36 artículos incluida en la descripción de YouTube, la estructura de capítulos, la portada y diagramas didácticos propios. La versión de YouTube muestra una portada estática o una imagen fija, por lo que no se insertan fotogramas repetidos de los participantes; al obtener localmente los metadatos de la versión de Bilibili con pantalla compartida que se menciona en la descripción, el servidor devolvió HTTP 412, y la extracción de la presentación de Feishu no obtuvo el contenido, de modo que esta versión se basa en las fuentes locales confiables.\par
\textbf{Página del proyecto}: \href{\repourl}{\nolinkurl{\repourl}}\par
\end{tcolorbox}
\end{titlepage}

\tableofcontents
\newpage
\section{Introducción: por qué entrar al mundo de la IA a través de los artículos}

Esta sección establece primero la actitud de estudio para este episodio. Xie Qingchi (谢青池) no presenta los artículos como investigador, sino como responsable de producto, autodidacta de largo plazo y observador de las fronteras de la tecnología, y hace público el recorrido que siguió durante más de un año para leer más de 200 artículos. Esta perspectiva es muy valiosa: no busca convertir cada artículo en una clase de matemáticas, sino preguntar qué frontera cambió cada artículo, qué paradigma abrió, quién lo heredó después y cómo se convirtió en la base de productos y de la industria.

Los 36 artículos de este episodio pueden condensarse en cuatro líneas principales. La primera son los paradigmas de modelos: de GPU, AlexNet, Attention y Transformer a MoE, CoT, LoRA, ReAct y Bitter Lesson. La segunda son la Infra y los datos: ZeRO (Zero Redundancy Optimizer, que reduce la redundancia del entrenamiento al fragmentar los estados del optimizador, los gradientes y los parámetros), Scaling Law, Chinchilla, LAION, RefinedWeb y MegaScale explican por qué los modelos pueden escalarse. La tercera son los modelos de lenguaje: Word2Vec, Google Translate, GPT, BERT, GPT-3, InstructGPT y Tulu 3 explican cómo los LLM se convirtieron en la puerta de entrada actual. La cuarta es la multimodalidad: DeepVideo, las redes de dos flujos, GAN, Diffusion, ViT, CLIP, Stable Diffusion y DiT explican cómo los modelos visuales y generativos confluyeron en la era del Transformer.

\begin{figure}[H]
\centering
\includegraphics[width=0.94\textwidth]{four-part-history-map.png}
\caption{Las cuatro líneas principales de los 36 artículos: los paradigmas de modelos, la Infra/los datos, los modelos de lenguaje y la multimodalidad constituyen en conjunto la historia de la evolución de la IA. Diagrama conceptual propio, elaborado a partir de la descripción del video y de la conversación de 00:19:35--03:56:38.}
\end{figure}

\begin{knowledgebox}{Lectura de la figura: no es una lista de artículos, sino un mapa de fronteras}
Los paradigmas de modelos determinan «qué estructuras pueden aprender»; la Infra y los datos determinan «hasta qué escala pueden crecer»; los modelos de lenguaje determinan «cómo las personas invocan la inteligencia con lenguaje natural»; la multimodalidad determina «cómo los modelos entran al mundo de la visión, el video y la generación». Solo al entrelazarse estas cuatro líneas se formó la industria de la IA actual.
\end{knowledgebox}

\begin{importantbox}{Tesis central de este episodio}
El valor de leer artículos no está en memorizar sus títulos, sino en entrar en contacto directo con el origen de los problemas. Comprender de verdad un artículo clave equivale a saber a qué frontera se enfrentaban los investigadores de su momento, por qué eligieron esa solución, qué parte amplificó después la época y qué parte quedó descartada por el hardware o los datos.
\end{importantbox}

\subsection{Aprender IA con IA: el ciclo de lectura de artículos}

El apartado anterior explicó por qué leer artículos; esta sección aborda cómo sostener esa lectura. Xie Qingchi cuenta que los obstáculos que encontró al principio incluían la base matemática, los artículos en inglés, la polisemia de los términos y la falta de una hoja de ruta. La solución no fue resistir a fuerza de voluntad, sino usar la IA para traducir, explicar, hacer preguntas de seguimiento y organizar, de modo que el artículo dejara de ser un PDF aislado y se convirtiera en un nodo de conocimiento que puede discutirse, repasarse y compartirse.

\begin{figure}[H]
\centering
\includegraphics[width=0.96\textwidth]{paper-reading-loop.png}
\caption{El ciclo de lectura de artículos para aprender IA con IA: la traducción, las preguntas, la hoja de ruta, el repaso y el compartir hacen que el artículo deje de ser solo un PDF. Diagrama conceptual propio, elaborado a partir de la conversación de 00:01:30--00:19:35.}
\end{figure}

\begin{knowledgebox}{Lectura de la figura: la IA no lee por ti, sino que te ayuda a levantar un andamiaje}
La IA puede ayudar a superar las barreras del inglés, de la terminología y de los conocimientos previos, pero no puede sustituir la capacidad de plantear problemas. Una buena forma de leer es encontrar primero el problema, después pedirle a la IA que ayude a traducir, explicar, comparar y cuestionar, y por último volver a situar el artículo en la hoja de ruta histórica.
\end{knowledgebox}

\subsection{Resumen de la sección}

El verdadero tema de EP117 es «cómo entrar al mundo de la tecnología». Los 36 artículos son solo el vehículo; el entrenamiento más profundo consiste en construir un mapa de problemas: saber qué conceptos permanecen estables a largo plazo, qué fronteras están cambiando, qué capacidades provienen del modelo y cuáles provienen de la Infra, los datos y el hardware.
\section{Cambio de paradigma en los modelos: de la GPU al Agent}

Esta sección cubre el primer eje: el paradigma de los modelos. La historia empieza con la primera GPU, en 1999. No es que la GPU sea un artículo de investigación, sino que cambió el límite de lo computable. Trabajos como Brook, CUDA, AlexNet, ResNet y Transformer muestran un hecho sencillo: que un paradigma algorítmico se vuelva el eje de una época depende muchas veces de que sepa aprovechar el hardware y los datos.

\begin{figure}[H]
\centering
\includegraphics[width=0.96\textwidth]{model-paradigm-timeline.png}
\caption{Línea de tiempo de los paradigmas de modelos: GPU, CNN, Attention, Transformer, MoE y Agent se relevan uno a otro. Diagrama conceptual propio, elaborado a partir de la conversación entre 00:19:35 y 01:52:58.}
\end{figure}

\begin{knowledgebox}{Lectura de la figura: los paradigmas no surgen de repente}
Antes de AlexNet ya existían la GPU y los datos de imágenes; antes de Transformer existían seq2seq, Attention, las redes residuales y el cómputo en paralelo; antes del Agent existían CoT, la llamada a herramientas y ReAct. Un cambio de paradigma suele ser la conexión repentina de muchos avances acumulados dentro de una ventana determinada de hardware y datos.
\end{knowledgebox}

\subsection{GPU, Brook, AlexNet y la lotería del hardware}

Esta sección examina primero cómo el hardware cambió el destino de los modelos. Brook for GPU representa la exploración temprana del cómputo de propósito general en GPU, mientras que AlexNet, en 2012, combinó la GPU con redes convolucionales profundas para superar ampliamente a todos en ImageNet. El ponente insiste en la «lotería del hardware»: si un método se ajusta precisamente al hardware dominante, se amplifican su velocidad de entrenamiento, su ecosistema de ingeniería y la inversión en investigación; en cambio, una idea puede ser correcta y aun así tener dificultades para volverse el eje de una época.

\begin{figure}[H]
\centering
\includegraphics[width=0.94\textwidth]{hardware-lottery.png}
\caption{Lotería del hardware: las arquitecturas que aprovechan el hardware dominante obtienen un rendimiento compuesto a largo plazo. Diagrama conceptual propio, elaborado a partir de la conversación entre 00:19:35 y 00:30:00, y entre 03:56:38 y 04:03:00.}
\end{figure}

\begin{knowledgebox}{Lectura de la figura: por qué AlexNet marca un parteaguas}
AlexNet no solo «usó una CNN»: también demostró que el aprendizaje profundo, la GPU, el conjunto de datos y las técnicas de ingeniería podían combinarse en una ventaja abrumadora. Derrotó al paradigma dominado por las características diseñadas a mano y anticipó la Bitter Lesson: a medida que el poder de cómputo sigue creciendo, los métodos de aprendizaje generales superan una y otra vez al diseño manual.
\end{knowledgebox}

\subsection{seq2seq, Attention y Transformer}

Lo anterior trató sobre imágenes; esta sección pasa a las secuencias. seq2seq hace que el modelo comprima la secuencia de entrada en un estado y luego la decodifique, mientras que Attention permite que, al decodificar, el modelo atienda directamente a distintas posiciones de la entrada, lo que mitiga los problemas de las oraciones largas y la pérdida de contexto. Transformer va más allá: sustituye el cuello de botella secuencial de las RNN por self-attention y cómputo en paralelo, de modo que puede modelar directamente las relaciones entre los token dentro de la secuencia.

\begin{figure}[H]
\centering
\includegraphics[width=0.96\textwidth]{transformer-origin-ladder.png}
\caption{Del modelado de secuencias a Transformer: seq2seq, Attention, las conexiones residuales y el cómputo en paralelo prepararon el terreno en conjunto. Diagrama conceptual propio, elaborado a partir de la conversación entre 00:34:02 y 01:05:24.}
\end{figure}

\begin{importantbox}{La intuición central de Attention}
Attention permite que el modelo, al procesar la posición actual, calcule directamente su relación con las demás posiciones. Frente a las RNN, que dependen de transmitir el estado oculto paso a paso, Attention acorta el recorrido de la información; frente a la convolución, que necesita varias capas para ampliar el campo receptivo, self-attention puede ver toda la ventana de contexto en una sola capa.
\end{importantbox}

\subsection{ResNet, destilación, AlphaGo Zero y MoE}

Lo anterior trató sobre cómo las arquitecturas procesan imágenes y secuencias; esta sección amplía el alcance a las técnicas de entrenamiento, la transferencia de capacidades y la economía del cómputo. El paradigma de los modelos no es solo la estructura: también incluye el entrenamiento y la transferencia de capacidades. ResNet usa conexiones residuales para resolver la dificultad de entrenar redes más profundas; la destilación plantea que un modelo pequeño puede aprender el conocimiento de un modelo grande; AlphaGo Zero mostró que el aprendizaje por refuerzo y el juego contra sí mismo pueden prescindir de las partidas humanas; y el MoE moderno reduce el costo de activación mediante el enrutamiento entre expertos, de modo que la capacidad de un modelo grande y su costo de cómputo ya no están ligados de forma completamente lineal.

\begin{knowledgebox}{Términos clave: hitos del paradigma de modelos}
\begin{tabularx}{\textwidth}{p{0.20\textwidth}p{0.32\textwidth}X}
\toprule
Hito & Problema que resuelve & Influencia posterior \\
\midrule
ResNet & Degradación y dificultad de entrenamiento de las redes profundas & Lo residual se volvió un componente general de las redes profundas y llegó también a Transformer. \\
Distillation & Cómo transferir las capacidades de un modelo grande a uno pequeño & Sustenta el despliegue en el dispositivo, los modelos estudiante y la compresión de modelos. \\
AlphaGo Zero & Aprender una estrategia sin depender de datos humanos & Las ideas del aprendizaje por refuerzo, el juego contra sí mismo y el test-time search se citan una y otra vez. \\
MoE & Aumentar la capacidad de parámetros controlando el cómputo de cada pasada & Modelos como DeepSeek convirtieron el enrutamiento entre expertos en una vía importante de los LLM modernos. \\
\bottomrule
\end{tabularx}
\end{knowledgebox}

\subsection{CoT, LoRA, ReAct y la transición hacia el Agent}

Esta sección pasa de los modelos base a la forma de usarlos. CoT hizo ver que la entrada del modelo influye en su salida, lo que impulsó la prompt engineering y la context engineering; LoRA generalizó la adaptación de parámetros a bajo costo; ReAct puso el reasoning y el acting en el mismo ciclo, lo que sentó las bases para que el Agent pasara de la teoría a la llamada a herramientas.

\begin{figure}[H]
\centering
\includegraphics[width=0.96\textwidth]{adaptation-agent-stack.png}
\caption{La pila que va de la capacidad a la acción: la destilación, LoRA, CoT y ReAct llevan al modelo hacia un sistema utilizable. Diagrama conceptual propio, elaborado a partir de la conversación entre 01:26:40 y 01:37:10.}
\end{figure}

\begin{knowledgebox}{Lectura de la figura: las capacidades del modelo deben «extraerse» y «recibirse»}
La destilación y LoRA resuelven la transferencia de capacidades y el costo de adaptación; CoT hace que el modelo despliegue su razonamiento de forma explícita, y ReAct conecta ese razonamiento con herramientas externas. El Agent no surgió de la nada: estos mecanismos fueron llevando poco a poco al modelo hacia un ciclo cerrado de acción.
\end{knowledgebox}

\subsection{Bitter Lesson: un baño de realidad para el largo plazo}

Todos los hitos anteriores insinúan una regularidad, y esta sección la expone por completo con The Bitter Lesson. Su tesis es que, a largo plazo, los métodos que dependen del cómputo general y de la búsqueda o el aprendizaje a gran escala suelen superar a los que codifican el conocimiento a mano. El ponente subraya que esto no significa que las características diseñadas a mano sean siempre inútiles, sino que las estructuras manuales suelen funcionar dentro de cierto orden de magnitud de cómputo; una vez que el cómputo y los datos superan un umbral, los métodos de aprendizaje generales vuelven a imponerse.

\begin{figure}[H]
\centering
\includegraphics[width=0.94\textwidth]{bitter-lesson-map.png}
\caption{Bitter Lesson: a largo plazo, el cómputo general y el aprendizaje a escala superan una y otra vez a las estructuras diseñadas a mano. Diagrama conceptual propio, elaborado a partir de la conversación entre 01:45:00 y 01:52:40.}
\end{figure}

\begin{warningbox}{Bitter Lesson no significa «no diseñar estructuras»}
El diseño de estructuras sigue siendo importante, sobre todo cuando hay pocos datos, poco cómputo y un objetivo claro. Pero si una estructura no puede hacer scale junto con el cómputo, los datos y el ecosistema de hardware, podría quedar superada por los métodos generales en la siguiente transición de paradigma.
\end{warningbox}

\subsection{Resumen de la sección}

El eje del paradigma de los modelos nos enseña que el progreso de la IA no es el triunfo de un solo artículo, sino la alineación conjunta del hardware, los datos, los métodos de entrenamiento, la arquitectura y la forma de uso. La GPU hizo viable el aprendizaje profundo, Attention hizo aprendibles las dependencias largas, Transformer aprovechó el hardware paralelo, y CoT/ReAct llevaron al modelo hacia la acción.
\section{Infra y datos: la base del aprendizaje a gran escala}

La sección anterior explicó «por qué cambia el modelo»; esta explica «por qué el modelo puede escalarse». Cuando crecen los parámetros del modelo, los datos y el contexto, el problema deja de ser solo algorítmico y pasa a abarcar la memoria de la GPU, la comunicación, el paralelismo, la calidad de los datos y la estabilidad del clúster. ZeRO es Zero Redundancy Optimizer, que reduce la redundancia del entrenamiento al particionar los estados del optimizador, los gradientes y los parámetros; Scaling Law, Chinchilla, LAION, RefinedWeb y MegaScale forman las demás líneas principales de esta parte.

\begin{figure}[H]
\centering
\includegraphics[width=0.86\textwidth]{infra-data-stack.png}
\caption{Pila de infra y datos: ZeRO, Scaling Law, los datos abiertos y los clústeres de decenas de miles de GPU escalan juntos el modelo. Diagrama conceptual propio, elaborado a partir de la conversación de 01:52:58--02:21:29.}
\end{figure}

\begin{knowledgebox}{Lectura de la figura: la base no es una tarea secundaria de trastienda}
Sin el paralelismo y la partición de estados de ZeRO, un modelo grande no cabe en varias GPU; aquí los estados incluyen el optimizer state, es decir, los estados del optimizador como \(m\) y \(v\) de Adam/AdamW. Sin Scaling Law ni Chinchilla, no se sabe cómo distribuir el presupuesto de entrenamiento; sin la ingeniería de datos de LAION o RefinedWeb, el modelo no tiene un corpus utilizable suficiente; sin la ingeniería de clústeres de MegaScale, entrenar con decenas de miles de GPU se queda en un eslogan.
\end{knowledgebox}

\subsection{ZeRO: particionar los estados}

Esta sección aborda primero un concepto de sistemas que es indispensable asimilar. La idea central de ZeRO es particionar los estados que, durante el entrenamiento, normalmente se replican en cada GPU, incluidos los estados del optimizador, los gradientes y los parámetros, para reducir la presión sobre la memoria de cada GPU. No es un algoritmo que haga más inteligente al modelo, pero hace posible entrenar modelos más grandes.

\begin{importantbox}{Qué es ZeRO: Zero Redundancy Optimizer}
ZeRO es la abreviatura de Zero Redundancy Optimizer. Su idea central es particionar (sharding) los estados redundantes del paralelismo de datos, de modo que cada GPU guarde solo una parte de los estados de entrenamiento en lugar de una copia completa. ZeRO-1 particiona los estados del optimizador, ZeRO-2 particiona además los gradientes y ZeRO-3 particiona además los parámetros. Aquí, el optimizer state suele referirse al momento de primer orden \(m\) y al momento de segundo orden \(v\) de Adam/AdamW, estados que aumentan considerablemente la memoria de la GPU necesaria para entrenar.
\end{importantbox}

\subsection{Scaling Law y Chinchilla: la guía para distribuir el presupuesto}

Scaling Law intenta describir cómo varía el desempeño del modelo en función del número de parámetros, la cantidad de datos y el cómputo. Chinchilla advierte además que, con un presupuesto de cómputo fijo, los parámetros del modelo y los tokens de entrenamiento deben guardar una proporción más adecuada; acumular parámetros sin datos suficientes puede no ser óptimo. Su importancia no está en ofrecer leyes eternas, sino en convertir el entrenamiento, antes decidido por intuición, en un problema predecible de distribución del presupuesto.

\begin{knowledgebox}{Fórmula intuitiva de Scaling Law}
La disminución de la pérdida puede entenderse de forma aproximada así:
\[
L(N, D, C) \downarrow \quad \text{as} \quad N,D,C \uparrow
\]
donde \(N\) representa el número de parámetros del modelo, \(D\) la cantidad de datos de entrenamiento y \(C\) el presupuesto de cómputo. Esta expresión no es una fórmula exacta, sino un recordatorio: los parámetros, los datos y el cómputo deben crecer de forma coordinada, sin apostar todo a una sola dimensión.
\end{knowledgebox}

\subsection{LAION, RefinedWeb y los datos abiertos}

La sección anterior trató cómo distribuir el presupuesto; esta pasa al combustible que hay detrás del presupuesto: los datos. LAION-5B refleja el esfuerzo de la comunidad de código abierto por romper el monopolio de unos pocos gigantes industriales; RefinedWeb muestra que los datos de internet, una vez limpiados, filtrados y combinados en proporciones adecuadas, pueden sostener modelos potentes. Más datos no siempre es mejor: el límite superior del modelo lo determinan en conjunto la procedencia, la calidad, el filtrado, la deduplicación, los derechos de autor y la reproducibilidad.

\begin{warningbox}{Los datos abiertos no son gratuitos}
Los datos abiertos reducen la barrera de entrada, pero también traen problemas de calidad, derechos de autor, sesgos y seguridad. La verdadera ingeniería de datos no consiste en descargar un archivo grande, sino en filtrar, deduplicar, etiquetar, evaluar y dar seguimiento de forma continua.
\end{warningbox}

\subsection{MegaScale: la lección de ingeniería del entrenamiento con decenas de miles de GPU}

Los datos determinan qué puede aprender el modelo, pero para suministrárselos realmente hace falta ingeniería de clústeres. En esta sección, MegaScale representa la realidad de ingeniería del entrenamiento a gran escala. En un clúster de entrenamiento, el costo real lo determinan la utilización de las GPU, la sobrecarga de comunicación, la recuperación ante fallas, la planificación, la topología de red y el ancho de banda de memoria. El ponente menciona MFU, es decir, Model FLOPs Utilization, que refleja la proporción de la capacidad de cómputo del hardware que el modelo aprovecha realmente; aun cuando acercarse al 50\% ya es excelente, eso significa que una gran parte del cómputo se pierde en comunicación y esperas.

\begin{knowledgebox}{Términos clave: infra y datos}
\begin{tabularx}{\textwidth}{p{0.22\textwidth}p{0.34\textwidth}X}
\toprule
Término & Problema que resuelve & Por qué importa \\
\midrule
ZeRO & Partición de los estados de entrenamiento & Permite que modelos más grandes quepan en la memoria de varias GPU. \\
Scaling Law & Predecir la relación entre escala y desempeño & Orienta el presupuesto de parámetros, datos y cómputo. \\
Chinchilla & Proporción entre datos y parámetros & Evita acumular parámetros con tokens de entrenamiento insuficientes. \\
LAION/RefinedWeb & Datos abiertos multimodales y web & Hacen que el entrenamiento de modelos no dependa por completo de gigantes cerrados. \\
MFU & Medir la utilización del hardware & Un MFU bajo hace que el mismo cómputo rinda menos avance de entrenamiento. \\
\bottomrule
\end{tabularx}
\end{knowledgebox}

\subsection{Resumen de la sección}

La línea de infra y datos muestra que los modelos grandes no crecen solo gracias a las fórmulas de los artículos. Requieren partición de la memoria, comunicación en paralelo, leyes de escala, ingeniería de datos, comunidad de código abierto y confiabilidad del clúster. El paradigma del modelo define la dirección; la infra y los datos determinan si esa dirección puede materializarse a escala.
\section{Modelos de lenguaje: de los vectores de palabras al posentrenamiento}

Las dos secciones anteriores explicaron los paradigmas de modelos y la base de entrenamiento; ahora volvemos a la puerta de entrada que hoy resulta más familiar para la mayoría: los modelos de lenguaje. Esta sección une en una sola línea «representación, preentrenamiento y posentrenamiento». La evolución de los modelos de lenguaje puede verse en tres etapas: primero, convertir las palabras en vectores; después, aprender las regularidades del lenguaje mediante el preentrenamiento; y, por último, usar instrucciones, preferencias y posentrenamiento de código abierto para que el modelo sea utilizable, controlable y capaz de dar servicio a los usuarios.

\begin{figure}[H]
\centering
\includegraphics[width=0.96\textwidth]{language-model-lineage.png}
\caption{Genealogía de los modelos de lenguaje: de los vectores de palabras al preentrenamiento y, después, a la alineación por instrucciones y al posentrenamiento de código abierto. Diagrama conceptual propio, elaborado a partir de la conversación de 02:21:29--03:08:08.}
\end{figure}

\begin{knowledgebox}{Lectura de la figura: el eje de los modelos de lenguaje va «de la representación a la conducta»}
Word2Vec resuelve la representación de las palabras; la traducción automática neuronal demostró que los modelos profundos pueden dar servicio en línea; GPT/BERT establecieron el paradigma del preentrenamiento; GPT-3 mostró la emergencia de capacidades con la escala; e InstructGPT y Tulu 3 llevaron los modelos hacia un uso práctico para las personas y hacia el posentrenamiento de código abierto.
\end{knowledgebox}

\subsection{Word2Vec y Google Translate: representación y despliegue}

Esta sección comienza con los antecedentes de los modelos de lenguaje, porque sin el aprendizaje de representaciones no habría existido el aprendizaje en contexto posterior. Word2Vec usa aprendizaje automático para convertir palabras en vectores, de modo que las relaciones semánticas entran en un espacio continuo. El despliegue de redes neuronales en Google Translate mostró, a su vez, que el aprendizaje profundo no solo funciona en los artículos académicos, sino que también puede sustituir a los métodos tradicionales en sistemas en línea a gran escala. Estos dos pasos respondieron, respectivamente, a «cómo representar el lenguaje» y a «cómo un modelo neuronal puede dar servicio a usuarios reales».

\begin{importantbox}{El significado de los vectores de palabras}
Los vectores de palabras proyectan palabras discretas en un espacio continuo, de modo que las palabras similares quedan más cerca en el espacio vectorial. Esto preparó el camino para las representaciones contextuales, el embedding del Transformer y la búsqueda semántica posteriores; sin embargo, sigue siendo una representación estática que no puede entender el significado de una palabra de forma dinámica según el contexto, como lo hace un LLM.
\end{importantbox}

\subsection{GPT-1, BERT, GPT-2 y GPT-3}

A partir de los vectores de palabras y la traducción en línea, los modelos de lenguaje entran de lleno en la era del preentrenamiento. GPT-1, BERT, GPT-2 y GPT-3, que se tratan en esta subsección, representan etapas distintas: GPT-1 propuso un nuevo paradigma de NLP basado en preentrenamiento no supervisado más fine-tuning supervisado; BERT, con masked language modeling bidireccional, se convirtió en «el rey de su época» en las tareas de comprensión; GPT-2 reforzó la ruta del preentrenamiento generativo e hizo que se empezara a considerar en serio la posibilidad de «prescindir del fine-tuning»; y GPT-3 combinó la Scaling Law, la apuesta organizacional y la ingeniería a gran escala, hasta convertirse en la señal más importante en la víspera de ChatGPT.

\begin{knowledgebox}{La bifurcación entre GPT y BERT}
\begin{tabularx}{\textwidth}{p{0.22\textwidth}p{0.34\textwidth}X}
\toprule
Ruta & Método de entrenamiento & Ventaja típica \\
\midrule
BERT & Predicción de palabras enmascaradas, contexto bidireccional & Destaca en tareas discriminativas como comprensión, clasificación y extracción. \\
GPT & next-token prediction autorregresiva & Gran capacidad de extensión a generación, completado, diálogo y tareas generales. \\
GPT-3 & Preentrenamiento autorregresivo a gran escala & El aprendizaje con pocos ejemplos, el aprendizaje en contexto y las capacidades emergentes pasan a ser el centro de atención. \\
\bottomrule
\end{tabularx}
\end{knowledgebox}

\subsection{InstructGPT y Tulu 3: llevar el modelo a la sociedad civilizada}

El significado central de InstructGPT es haber llevado al modelo de «saber continuar texto de internet» a «poder responder según la intención humana». Detrás de esto están métodos de posentrenamiento como los datos de instrucciones, los datos de preferencias y el RLHF. Tulu 3, por su parte, representa el intento de la comunidad de código abierto de sistematizar y hacer transparente el proceso de posentrenamiento, para que las capacidades de los modelos no queden solo en manos de unos pocos laboratorios de código cerrado.

\begin{warningbox}{El posentrenamiento no es un simple complemento}
El preentrenamiento le da al modelo conocimiento del mundo y capacidad lingüística; el posentrenamiento determina cómo responde a las personas, cómo sigue instrucciones, cómo rechaza solicitudes peligrosas, cómo produce salidas con el formato indicado y cómo completa tareas. Sin posentrenamiento, incluso un modelo potente puede tener dificultades para convertirse en un buen producto.
\end{warningbox}

\subsection{Resumen de la sección}

La línea de los modelos de lenguaje llevó a la IA del aprendizaje de representaciones a una puerta de entrada de interacción general. De Word2Vec a GPT-3 hay un crecimiento de capacidades; de InstructGPT a Tulu 3, un crecimiento en facilidad de uso, alineación y posentrenamiento de código abierto. Los productos de LLM actuales se construyen en el punto donde se cruzan estas dos curvas de crecimiento.
\section{Multimodalidad: de la comprensión de video a la fusión de modelos generativos}

Esta sección aborda la cuarta línea principal: la multimodalidad. El desarrollo multimodal no consiste simplemente en «agregarle imágenes a un modelo de lenguaje», sino en la convergencia gradual de la comprensión visual, el modelado de video, los modelos generativos, el Transformer visual, la alineación imagen-texto y el Transformer de difusión. La palabra central aquí es fusión: las distintas modalidades terminan por influirse mutuamente dentro de una representación unificada, una generación unificada o un marco de razonamiento unificado.

\begin{figure}[H]
\centering
\includegraphics[width=0.96\textwidth]{multimodal-lineage.png}
\caption{Genealogía de los modelos multimodales: convergencia de video, GAN, Diffusion, ViT, CLIP, Stable Diffusion y DiT. Diagrama conceptual propio, elaborado a partir de la conversación entre 03:08:08 y 03:56:38.}
\end{figure}

\begin{knowledgebox}{Lectura de la figura: la multimodalidad no es una línea recta}
La comprensión de video, la generación de imágenes, el Transformer visual, la alineación imagen-texto y los modelos de difusión partieron de problemas distintos. Convergieron después porque los modelos necesitan comprender el mundo y generarlo a la vez, y conectar la capacidad visual con la intención del usuario por medio del lenguaje.
\end{knowledgebox}

\subsection{DeepVideo, redes de dos flujos y la comprensión temprana de video}

DeepVideo y las redes de dos flujos representan los primeros intentos del aprendizaje profundo por entrar al ámbito del video. Una vez que las imágenes pudieron procesarse con CNN, los investigadores se preguntaron con naturalidad si el video también podía tratarse así. La dificultad es que el video añade la dimensión temporal: el modelo no solo debe reconocer el contenido de cada cuadro, sino también comprender acciones, movimiento y eventos.

\begin{importantbox}{Los problemas que el video añade frente a la imagen}
El modelado de video añade al menos tres problemas: la dependencia temporal, la representación del movimiento y los eventos de largo alcance. Una imagen aislada responde «qué es esto»; el video además debe responder «cómo cambia», «por qué cambia así» y «qué ocurrirá después».
\end{importantbox}

\subsection{GAN, Diffusion y DDPM}

La línea de generación de imágenes estuvo dominada primero por las GAN. Las GAN generan muestras nítidas, pero su entrenamiento es inestable; los VAE se entrenan de forma estable, pero sus imágenes tienden a salir borrosas. Diffusion creció en sus inicios a la sombra de las GAN, y DDPM devolvió a los modelos de difusión al centro del escenario de la generación de imágenes. La ventaja de los modelos de difusión es que su entrenamiento es estable, la calidad de generación es buena y se combinan bien con el control condicional y la guía por texto.

\begin{knowledgebox}{Ventajas y desventajas de GAN frente a Diffusion}
\begin{tabularx}{\textwidth}{p{0.22\textwidth}p{0.34\textwidth}X}
\toprule
Familia de modelos & Ventajas & Dificultades \\
\midrule
GAN & Muestras nítidas, generación rápida & Entrenamiento inestable, alto riesgo de colapso de modos. \\
VAE & Entrenamiento estable, interpretación como modelado probabilístico & Las muestras tienden a salir borrosas. \\
Diffusion & Estable, de alta calidad, fácil de controlar condicionalmente & Costo de muestreo elevado; requiere optimización de ingeniería. \\
\bottomrule
\end{tabularx}
\end{knowledgebox}

\subsection{ViT, CLIP, Stable Diffusion y DiT}

La subsección anterior mostró cómo los modelos generativos avanzaron hacia la difusión; esta examina cómo la visión y el lenguaje convergen en una arquitectura unificada. ViT divide la imagen en patches y procesa los tokens visuales con un Transformer; CLIP coloca imágenes y textos en un mismo espacio semántico y se convirtió en un pilar de la generación de imágenes a partir de texto y de la búsqueda multimodal; Stable Diffusion combinó los modelos de difusión, el espacio latente y el ecosistema de código abierto, e impulsó la popularización de la generación de imágenes; DiT, por su parte, fusiona más a fondo Diffusion y Transformer, y representa la expectativa de una futura arquitectura unificada.

\begin{importantbox}{Por qué CLIP es fundamental}
La importancia de CLIP reside en la alineación imagen-texto. Permite que el modelo sepa si una imagen y un texto coinciden semánticamente, lo que sustenta la generación de imágenes a partir de texto, la búsqueda de imágenes, el filtrado de datos y el aprendizaje de representaciones multimodales. Sin alineación imagen-texto, a un modelo generativo le resulta muy difícil entender de manera estable las indicaciones del usuario.
\end{importantbox}

\subsection{Resumen de la sección}

La línea multimodal muestra que la visión, el lenguaje y la generación se están acercando. Las GAN establecieron el punto de partida de la generación, Diffusion resolvió la calidad y la estabilidad, ViT llevó la visión al Transformer, CLIP tendió el puente entre imagen y texto, y DiT terminó de fusionar la generación con el Transformer.

\section{Mapa de aprendizaje de código abierto y fronteras futuras}

Las cuatro secciones anteriores cubrieron la línea principal de artículos; esta sección vuelve a la conversación final. Los temas clave que discuten Zhang Xiaojun (张小珺) y Xie Qingchi (谢青池) incluyen: la arquitectura debe apoyarse firmemente en el hardware; la frontera tecnológica actual aún no toca techo; OpenAI podría pasar de ser un lab a convertirse en una súper app o un sistema operativo; Google conserva una base sólida de talento, ingeniería e Infra. Lo más valioso que conviene llevarse de aquí es la conciencia de las fronteras: la IA está todavía en una etapa muy temprana y, aunque muchas capacidades parecen asombrosas, si se compara con la historia de la computación personal, quizá todavía estemos en la era de las minicomputadoras o incluso antes de la computación personal.

\begin{figure}[H]
\centering
\includegraphics[width=0.96\textwidth]{open-source-learning-map.png}
\caption{Mapa de aprendizaje de código abierto: artículos, código, datos, modelos, PPT y comunidad forman una infraestructura de aprendizaje. Diagrama conceptual propio, elaborado a partir de la conversación de 03:56:38--04:22:20.}
\end{figure}

\begin{knowledgebox}{Lectura de la figura: el código abierto no es solo código}
Este episodio forma parte, en sí mismo, del aprendizaje de código abierto: la ruta de artículos, la explicación en video, las PPT, los conjuntos de datos, el código y la discusión de la comunidad reducen juntos la barrera de entrada. El valor del código abierto no está solo en reproducir resultados, sino también en permitir que más personas vean el origen de los problemas y su contexto histórico.
\end{knowledgebox}

\subsection{Recomendaciones para dos tipos de lectores}

Esta sección concreta las recomendaciones finales. Para quien observa el mundo de la IA desde fuera, lo más importante es construir primero un mapa: usar IA para traducir artículos, leer reseñas y mapas de ruta, y entender los conceptos que no cambiarán en tres a cinco años, sin hundirse de entrada en los detalles de las fórmulas. Para quien ya lleva años trabajando dentro del sistema, lo importante es volver al origen de los problemas: por qué surgió este artículo, cuáles son las fronteras del hardware y de los datos, y si los productos actuales malinterpretan el mecanismo del artículo.

\begin{figure}[H]
\centering
\includegraphics[width=0.94\textwidth]{two-reader-advice.png}
\caption{Recomendaciones para dos tipos de lectores: quien observa desde fuera construye primero un mapa; quien trabaja dentro del sistema debe cuestionar las fronteras. Diagrama conceptual propio, elaborado a partir de la conversación de 03:56:38--04:22:20.}
\end{figure}

\begin{knowledgebox}{Lectura de la figura: cada tipo de lector tiene prioridades de aprendizaje distintas}
Al lector externo lo que más le falta es un punto de entrada y confianza, por eso conviene que lea primero los mapas de ruta; el lector interno tiende a caer en la optimización local, por eso debe volver al origen de los problemas, a las restricciones de hardware y al contexto histórico. Ambos necesitan los artículos, pero los leen de manera distinta.
\end{knowledgebox}

\subsection{Hardware, súper app y sistema operativo}

La sección anterior dio recomendaciones de aprendizaje a los lectores; esta sección las sitúa en la competencia de la industria: se leen artículos para ver con claridad las fronteras de las plataformas. La discusión final sobre OpenAI y Google puede ubicarse en un marco de competencia entre plataformas. Google tiene una base de talento, ingeniería e Infra; si OpenAI fuera solo un lab, le resultaría difícil ponerse a la par de Google; si se convierte en una súper app o un sistema operativo, la estructura de la competencia cambia. Este juicio se corresponde con el Agent OS de EP118, la AI War de EP127 y la convergencia de las fronteras de los Agent de EP139: las empresas de modelos terminarán disputándose los puntos de entrada, las herramientas, el ecosistema y la posición de sistema operativo.

\begin{warningbox}{No confundir el liderazgo en modelos con la victoria como plataforma}
La capacidad del modelo es una condición necesaria, pero una plataforma también requiere puntos de entrada de producto, un ecosistema de desarrolladores, llamadas a herramientas, retorno de datos, un modelo de negocio y ejecución organizacional. La diferencia entre Google y OpenAI no está solo en los benchmark, sino en los recursos de base y en la trayectoria como plataforma.
\end{warningbox}

\subsection{Resumen de la sección}

El cierre de EP117 nos recuerda que la frontera tecnológica está lejos de agotarse; lo que de verdad hay que aprender son los problemas y las fronteras, no las conclusiones de un solo artículo. Un mapa de aprendizaje de código abierto puede reducir la barrera de entrada, pero la comprensión a largo plazo sigue surgiendo de la lectura constante, la revisión retrospectiva y el hábito de situar los artículos en su historia.
\section{Síntesis y ampliación}

Esta sección condensa el episodio completo en un marco de estudio. Primero, los paradigmas de modelos no son artículos aislados: son el resultado de que el hardware, los datos, la arquitectura y las técnicas de entrenamiento se alineen. Segundo, la Infra y los datos amplifican la capacidad de los modelos. Sin fragmentación de la memoria de GPU, sin eficiencia de clúster y sin datos de alta calidad, la Scaling Law se queda en una regularidad sobre el papel. Tercero, el desarrollo de los modelos de lenguaje avanza de la representación al preentrenamiento y después al posentrenamiento y la alineación. Cuarto, el desarrollo multimodal es una convergencia gradual de la comprensión visual, la generación de imágenes, la alineación entre imagen y texto y la adopción del Transformer. Quinto, lo esencial al leer artículos no es terminarlos, sino desarrollar conciencia de sus límites.

\begin{importantbox}{Ubicar EP117 en la serie de IA de Zhang Xiaojun (张小珺)}
EP119 trata la arqueología de la arquitectura Attention. EP118 explica cómo entiende un CEO la VLA, el Agent OS y las plataformas. EP117 lleva estas preguntas más atrás en la historia: por qué ganó el Transformer, por qué la Infra y los datos se volvieron la línea principal y por qué los modelos de lenguaje y los multimodales terminaron por converger.
\end{importantbox}

\subsection{Tabla condensada de los 36 artículos}

\begin{longtable}{p{0.22\textwidth}p{0.30\textwidth}p{0.38\textwidth}}
\toprule
Línea principal & Artículo/hito & Importancia didáctica \\
\midrule
Paradigmas de modelos & Brook, AlexNet, ResNet & La GPU, el aprendizaje profundo y las redes residuales cambiaron juntos el límite de lo que se puede entrenar. \\
Paradigmas de modelos & seq2seq, Attention, Transformer & El modelado de secuencias pasa de comprimir el estado a modelar directamente las relaciones entre tokens. \\
Paradigmas de modelos & AlphaGo Zero, MoE, CoT, LoRA, ReAct & El aprendizaje por refuerzo, el enrutamiento entre expertos, las indicaciones para razonar, la adaptación de bajo costo y los Agents forman las capacidades de los sistemas posteriores. \\
Paradigmas de modelos & The Bitter Lesson & A largo plazo, el cómputo general y el aprendizaje a escala superan una y otra vez a las estructuras diseñadas a mano. \\
Infra/datos & ZeRO, MegaScale & La fragmentación y la ingeniería de clústeres determinan si un modelo grande puede entrenarse. \\
Infra/datos & Scaling Law, Chinchilla & Los parámetros, los datos y el presupuesto de cómputo deben guardar proporción. \\
Infra/datos & LAION-5B, RefinedWeb & Los datos abiertos y la ingeniería de datos web determinan el límite del corpus entrenable. \\
Modelos de lenguaje & Word2Vec, Google Translate & La representación de palabras y el despliegue en línea de modelos neuronales son la prehistoria de los modelos de lenguaje. \\
Modelos de lenguaje & GPT-1, BERT, GPT-2, GPT-3 & El paradigma de preentrenamiento pasa de las tareas de comprensión a la generación y a la emergencia de capacidades por escala. \\
Modelos de lenguaje & InstructGPT, Tulu 3 & El posentrenamiento, la alineación con instrucciones y la reproducción abierta hacen que los modelos sean utilizables. \\
Multimodal & DeepVideo, redes de dos flujos & La comprensión de video lleva el aprendizaje profundo de la imagen estática a la dimensión temporal. \\
Multimodal & GAN, Diffusion, DDPM & Los modelos generativos pasan del entrenamiento adversarial a la generación estable por difusión. \\
Multimodal & ViT, CLIP, Stable Diffusion, DiT & El Transformer visual, la alineación entre imagen y texto y el Transformer de difusión convergen. \\
\bottomrule
\end{longtable}

\subsection{Lecturas adicionales}

\begin{itemize}
\item Si le interesan los paradigmas de modelos, puede contrastar este episodio con la revisión de la arquitectura Attention de EP119, para entender por qué, después del Transformer, todavía se sigue «puliendo la Attention».
\item Si le interesa la IA del mundo físico, puede contrastarlo con EP118, EP120 y EP121, para reunir la VLA, los modelos del mundo y los datos de robots en un mismo mapa.
\item Si le interesan los Agents, puede contrastarlo con la historia técnica de los Agents de EP139 y con el paradigma de posentrenamiento/Agent de EP138, para entender cómo CoT/ReAct se convierten en sistemas reales.
\end{itemize}

\end{document}
